\chapter{Kesimpulan dan Saran}
\label{chap:kesimpulan-saran}

\section{Kesimpulan}
\label{sec:kesimpulan}

Berdasarkan hasil perancangan, implementasi, dan pengujian yang telah dilakukan, sistem kehadiran berbasis pengenalan wajah menggunakan konsep \textit{Federated Learning} pada lingkungan \textit{Edge Computing} telah berhasil dibangun. Sistem ini diimplementasikan menggunakan model \textit{MobileFaceNet} di dalam \textit{container Docker} pada perangkat Jetson Nano (\textit{Client} 1) dan Raspberry Pi (\textit{Client} 2), dengan \textit{server} agregator terpusat. Melalui sistem ini, keamanan privasi data biometrik mahasiswa terjaga karena seluruh citra wajah mentah disimpan secara lokal di masing-masing perangkat \textit{client}, tanpa perlu dikirim ke \textit{server}.

Adapun hasil perbandingan performa model antara metode \textit{Federated Learning} dan \textit{Centralized Learning} adalah sebagai berikut:
\begin{enumerate}[itemsep=1pt]
    \item Kecerdasan model global pada \textit{Federated Learning} mampu menyeimbangi \textit{Centralized Learning} dengan hasil akurasi sebesar 98,16\% dan \textit{validation loss} 0,3599, sedangkan \textit{Centralized Learning} menghasilkan akurasi sebesar 96,93\% dan \textit{validation loss} 0,5457.
    \item Durasi pelatihan keseluruhan pada \textit{Federated Learning} berjalan jauh lebih lambat yaitu selama 5034,13 detik dibandingkan \textit{Centralized Learning} yang hanya membutuhkan waktu 439,99 detik, akibat adanya waktu tunggu sinkronisasi parameter antar-perangkat.
    \item Kecepatan pelatihan lokal pada tingkat \textit{client} dipengaruhi oleh spesifikasi RAM, di mana Raspberry Pi (RAM 1 GB) membutuhkan waktu latih rata-rata 334 hingga 384 detik per ronde, sedangkan Jetson Nano (RAM 4 GB) jauh lebih cepat dengan durasi 113 hingga 114 detik per ronde.
    \item Penggunaan \textit{bandwidth} jaringan pada \textit{Federated Learning} jauh lebih efisien sebesar 159,46 MB dibandingkan \textit{Centralized Learning} yang menghabiskan 1152,34 MB karena harus mengunggah seluruh citra mentah ke \textit{server}.
    \item Konsumsi energi listrik pada \textit{Federated Learning} lebih besar mencapai 0,089481 kWh akibat durasi pelatihan yang lama pada seluruh perangkat, sedangkan \textit{Centralized Learning} lebih hemat yaitu sebesar 0,012222 kWh karena pemrosesan terpusat selesai secara cepat di tingkat \textit{server}.
    \item Performa inferensi secara nyata (\textit{real-life}) pada kedua metode menghasilkan tingkat keberhasilan fungsionalitas 100\% pada jarak pemindaian 0,5 meter hingga 3 meter pada kondisi cahaya cukup dan minim dengan tingkat keberhasilan klasifikasi stabil di atas 90\% pada ambang batas kepercayaan 0,7. Latensi inferensi lokal pada Jetson Nano berjalan cepat ($\sim$700 ms) sedangkan pada Raspberry Pi berjalan lambat ($\sim$3100 ms).
\end{enumerate}

\section{Saran}
\label{sec:saran}

Berdasarkan keterbatasan sistem serta hasil evaluasi yang diperoleh sepanjang pelaksanaan penelitian ini, beberapa saran yang dapat diajukan untuk pengembangan penelitian selanjutnya adalah sebagai berikut:
\begin{enumerate}[itemsep=1pt]
    \item Menggunakan arsitektur model pengenalan wajah alternatif yang lebih modern selain \textit{MobileFaceNet} guna menganalisis perbandingan akurasi dan latensi inferensinya.
    \item Menerapkan metode deteksi keaktifan wajah (\textit{anti-spoofing}) untuk memperkuat keamanan sistem dari ancaman pemalsuan identitas menggunakan foto atau video.
    \item Menggunakan perangkat keras \textit{edge} dengan spesifikasi komputasi dan kapasitas RAM yang lebih tinggi (seperti \textit{Jetson Orin Nano} atau perangkat keras sejenis dengan kapasitas RAM di atas 4 GB) agar proses pelatihan terdistribusi dapat berjalan lebih cepat serta mampu menampung skala data latih yang lebih besar (lebih dari 30 subjek).
    \item Mengekspansi skala pengujian dengan melibatkan lebih dari 30 subjek mahasiswa untuk mengevaluasi keandalan serta skalabilitas model global hasil agregasi secara lebih luas.
\end{enumerate}
